\chapter{Linear algebra beyond the first course}\label{ch01}
% Short running heads for long sentence-style section titles. The original
% \sectionmark is saved here and restored at the end of the chapter.
\let\ompSavedSectionmark\sectionmark
\providecommand{\ompnexthead}[1]{\def\ompheadtext{#1}%
  \renewcommand{\sectionmark}[1]{\markright{\MakeUppercase{\thesection.\ \ompheadtext}}}}
\providecommand{\Kc}{K_c}

The first course in linear algebra ends with eigenvalues of small symmetric
matrices, the singular value decomposition and the projection onto a line.
Primer~L of \emph{Data Mining as Observation} covers that ground, and its
chapter~0 introduces the read operator $\PC$ and the observer distortion
$\dO=\tr(\PC\Sd)$. This chapter starts there. Its subject is the algebra of
\emph{two} quadratic forms at once, because nearly every object in the
programme is one quadratic form measured against another. An error covariance
is read through a consumer's operator. A source covariance sets the scale
against which a reader's weights are compared. A budget enters a pencil of two
matrices whose largest generalized eigenvalue decides a rate.

Each section introduces one tool, states it as a numbered equation or theorem,
and works an example small enough to check by hand. Every number in those
examples is recomputed by \texttt{checks/ch01\_examples.py}.

% ---------------------------------------------------------------------------
\ompnexthead{Inner products and projectors}
\section{A positive definite matrix defines an inner product, and an orthogonal projector is a symmetric idempotent}\label{sec:ch01:proj}

The dot product $x\T y$ of Primer~L, section~L.2, treats every coordinate
alike. Many readers do not. A consumer that is twice as sensitive to the first
coordinate as to the second measures length with a weight. Any symmetric
positive definite matrix $G$ supplies such a weight, and the result has every
property the reader uses an inner product for.

\begin{definition}[Weighted inner product]\label{def:ch01:ginner}
For a symmetric matrix $G\pd 0$, the $G$-inner product and the norm it induces are
\begin{equation}\label{eq:ch01:ginner}
\inner{x}{y}_G = x\T G\,y,\qquad \norm{x}_G=\sqrt{x\T G\,x}.
\end{equation}
\end{definition}

The inner product is symmetric because $G$ is, it is linear in each argument,
and $\norm{x}_G>0$ for $x\neq 0$ because $G$ is positive definite. When $G$ is
only positive semidefinite the same formula gives a \emph{seminorm}. It assigns
length zero to every vector in $\ker G$. That case is the normal one for a read
operator, whose kernel holds the directions a consumer never reads, and it is
why the programme writes $\PC\psd 0$ rather than $\PC\pd 0$.

Primer~L, section~L.5, builds the projection onto one unit vector as $uu\T$. The
general object is a matrix that does nothing the second time it is applied.

\begin{definition}[Projector]\label{def:ch01:projector}
A square matrix $\Pi$ is a \emph{projector} when it is idempotent, and an
\emph{orthogonal projector} when it is also symmetric,
\begin{equation}\label{eq:ch01:projector}
\Pi^2=\Pi \quad\text{(projector)},\qquad \Pi^2=\Pi=\Pi\T\quad\text{(orthogonal projector)}.
\end{equation}
\end{definition}

A projector splits every vector into a part it keeps and a part it discards,
$x=\Pi x+(I-\Pi)x$. The kept part lies in $\range\Pi$ and the discarded part
lies in $\ker\Pi$, since $\Pi(I-\Pi)x=(\Pi-\Pi^2)x=0$. Symmetry is exactly the
condition that makes the two parts perpendicular,
\begin{equation}\label{eq:ch01:split}
(\Pi x)\T(I-\Pi)y = x\T(\Pi\T-\Pi\T\Pi)y = 0\ \text{for all }x,y
\quad\Longleftrightarrow\quad \Pi\T=\Pi\T\Pi \quad\Longleftrightarrow\quad \Pi=\Pi\T .
\end{equation}
The last equivalence holds because $\Pi\T\Pi$ is symmetric, so $\Pi\T=\Pi\T\Pi$
forces $\Pi\T$ to be symmetric, and transposing $\Pi\T=\Pi\T\Pi$ then gives
$\Pi=\Pi\T\Pi=\Pi\T$. An idempotent that is not symmetric is an \emph{oblique}
projector. It still splits $x$ into a range part and a kernel part, but the two
parts are not at right angles.

If the columns of $Q\in\R^{d\times k}$ are an orthonormal basis of a subspace
$\mathcal S$, the orthogonal projector onto $\mathcal S$ is $QQ\T$. If the
columns of $A$ are any basis, orthonormal or not, the same projector is
\begin{equation}\label{eq:ch01:projbasis}
\Pi_{\mathcal S}=QQ\T=\sum_{i=1}^k q_iq_i\T = A(A\T A)^{-1}A\T,\qquad
\tr\Pi_{\mathcal S}=\rank\Pi_{\mathcal S}=k .
\end{equation}
The trace equals the rank because the eigenvalues of an orthogonal projector
are $1$ on $\mathcal S$ and $0$ on its complement, and the trace is the sum of
the eigenvalues.

Replacing the dot product by the $G$-inner product changes which point of
$\mathcal S$ is nearest. Minimizing $\norm{x-Ac}_G^2$ over $c$ gives the normal
equations $A\T G(x-Ac)=0$, and so the \emph{$G$-orthogonal projector}
\begin{equation}\label{eq:ch01:gproj}
\Pi_G = A(A\T GA)^{-1}A\T G,\qquad \Pi_G^2=\Pi_G,\qquad G\,\Pi_G=\Pi_G\T G .
\end{equation}
It is idempotent but in general not symmetric. It is orthogonal in the
$G$-inner product, which is what the identity $G\Pi_G=\Pi_G\T G$ says. Read
in the ordinary dot product, it is oblique unless it happens to be symmetric,
as it is when $G=cI$ or when $\mathcal S$ is spanned by eigenvectors of $G$. Because the residual is
$G$-orthogonal to the subspace, Pythagoras holds in the $G$-norm,
\begin{equation}\label{eq:ch01:gpyth}
\norm{x}_G^2=\norm{\Pi_Gx}_G^2+\norm{x-\Pi_Gx}_G^2 .
\end{equation}
The split of a weighted error into a part the reader sees in a subspace and a
part it sees outside is this identity.

\begin{example}[Two projections of the same point]\label{ex:ch01:gproj}
Take the line spanned by $u=(1,1)\T$, the point $x=(1,-1)\T$ and the weight
$G=\diag(2,1)$. The ordinary projector is $uu\T/2$, with every entry
$\tfrac12$, and it sends $x$ to $(0,0)\T$ because $x\perp u$. The
$G$-projector is $\Pi_G=u(u\T Gu)^{-1}u\T G$. Here $u\T Gu=3$ and
$u\T G=(2,1)$, so
\[
\Pi_G=\tfrac13\begin{pmatrix}2&1\\2&1\end{pmatrix},\qquad
\Pi_G x=\tfrac13\,(1,1)\T .
\]
Squaring $\Pi_G$ gives back $\Pi_G$, since $\tfrac19(4+2)=\tfrac23$ in the
first entry and the rest follow the same way, and $\Pi_G$ is visibly not
symmetric. The residual is $x-\Pi_Gx=(\tfrac23,-\tfrac43)\T$. Its $G$-inner
product with $u$ is $\tfrac23\cdot2-\tfrac43\cdot1=0$, so it is $G$-orthogonal
to the line, while its dot product with $u$ is $-\tfrac23\neq0$. The squared
$G$-distance from $x$ to the line is $2\cdot\tfrac49+\tfrac{16}9=\tfrac83$,
smaller than the $G$-distance $\norm{x}_G^2=3$ to the Euclidean foot $(0,0)\T$.
The kept part has $\norm{\Pi_Gx}_G^2=(2+1)/9=\tfrac13$, and
$\tfrac13+\tfrac83=3$, as \cref{eq:ch01:gpyth} requires.
\Cref{fig:ch01:gproj} shows the two feet.

For a subspace of dimension two in $\R^3$, take $q_1=(1,1,0)\T/\sqrt2$ and
$q_2=(0,0,1)\T$. Then
\[
\Pi=q_1q_1\T+q_2q_2\T=\begin{pmatrix}\tfrac12&\tfrac12&0\\ \tfrac12&\tfrac12&0\\ 0&0&1\end{pmatrix},
\qquad \Pi\,(3,1,2)\T=(2,2,2)\T .
\]
The residual $(1,-1,0)\T$ spans $\ker\Pi$, and $\tr\Pi=2=\rank\Pi$.
\Cref{fig:ch01:plane} draws the split.
\end{example}

\begin{figure}[tbp]
\centering
\begin{tikzpicture}[scale=1.35]
\draw[guide,->] (-1.3,0)--(2.6,0) node[lbl,right] {$x_1$};
\draw[guide,->] (0,-2.9)--(0,1.4) node[lbl,above] {$x_2$};
\draw[thick] (-1.2,-1.2)--(1.3,1.3) node[lbl,above right] {$\operatorname{span}(u)$};
\draw[thick,pblue] (1,-1) circle ({sqrt(2)});
\draw[thick,porange] (1,-1) ellipse ({sqrt(4/3)} and {sqrt(8/3)});
\fill (1,-1) circle (1.4pt) node[lbl,below right] {$x$};
\fill[pblue] (0,0) circle (1.8pt);
\node[lbl,pblue,left] at (-0.05,0.14) {$\Pi x$};
\fill[porange] (1/3,1/3) circle (1.8pt);
\node[lbl,porange,above left] at (1/3,1/3) {$\Pi_G x$};
\end{tikzpicture}
\caption{The Euclidean and the $G$-orthogonal projection of $x=(1,-1)\T$ onto the
line spanned by $(1,1)\T$, with $G=\diag(2,1)$. The circle is the smallest
Euclidean ball around $x$ that touches the line, and it touches at the origin.
The ellipse is the smallest $G$-ball, and it touches at $(\tfrac13,\tfrac13)$.}
\label{fig:ch01:gproj}
\end{figure}

\begin{figure}[tbp]
\centering
\begin{tikzpicture}[x={(-0.55cm,-0.42cm)},y={(1.1cm,0cm)},z={(0cm,1.1cm)},scale=1.1]
\fill[plight] (0,0,0)--(3.2,3.2,0)--(3.2,3.2,3)--(0,0,3)--cycle;
\draw[pgray] (0,0,0)--(3.2,3.2,0)--(3.2,3.2,3)--(0,0,3)--cycle;
\draw[guide,->] (0,0,0)--(3.8,0,0) node[lbl,below left] {$x_1$};
\draw[guide,->] (0,0,0)--(0,3.8,0) node[lbl,right] {$x_2$};
\draw[guide,->] (0,0,0)--(0,0,3.5) node[lbl,above] {$x_3$};
\draw[vec,pblue] (0,0,0)--(3,1,2) node[lbl,left] {$x=(3,1,2)$};
\draw[vec,pgreen] (0,0,0)--(2,2,2) node[lbl,above right] {$\Pi x=(2,2,2)$};
\draw[vec,pred] (2,2,2)--(3,1,2);
\node[lbl,pred] at (2.5,1.5,2.75) {$x-\Pi x$};
\node[lbl,pgray,right] at (0,0,3.15) {$\operatorname{span}(q_1,q_2)$};
\end{tikzpicture}
\caption{The orthogonal projector onto the plane spanned by $q_1$ and $q_2$ keeps
$(2,2,2)\T$, which lies in the plane, and discards $(1,-1,0)\T$, which is
perpendicular to it. The three vectors form a right triangle, so the squared
lengths add, $14=12+2$.}
\label{fig:ch01:plane}
\end{figure}

\inproject{The read operator is a weighted sum of projectors onto the directions a
consumer reads. In \nolinkurl{geometric-observation/paper/observation-theory.tex},
appendix ``The price of a misidentified observer'', the theorem ``Omission floor''
is stated with $\Pi$ the orthogonal projector onto a kernel, and the paragraph
after it uses the overlap $\tfrac1r\tr(\Pi_R\Pi_{\hat R})$ of two projectors to
measure how well a probe recovered a read subspace.}

\buildson{Primer~L, sections~L.4 and L.5 (kernel, rank, the projection $uu\T$).
Chapter~0 of \emph{Data Mining as Observation}, section~0.2.}
\usedin{\Cref{sec:ch01:whiten,sec:ch01:pencil}. \Cref{ch05}, where a
positive definite $G$ that varies from point to point is a Riemannian metric.
\Cref{ch06}, read subspaces and their overlaps.}

% ---------------------------------------------------------------------------
\ompnexthead{Spectral theorem and Loewner order}
\section{A symmetric matrix is a weighted sum of orthogonal projectors, and the Loewner order compares them}\label{sec:ch01:spectral}

Primer~L, section~L.7, writes a symmetric matrix as $V\Lambda V\T$. Grouping
the eigenvectors that share an eigenvalue turns that into a statement about
projectors, and in that form it handles repeated eigenvalues without fuss.

\begin{theorem}[Spectral theorem]\label{thm:ch01:spectral}
Let $A\in\R^{d\times d}$ be symmetric with distinct eigenvalues
$\lambda_1>\dots>\lambda_m$. Then
\begin{equation}\label{eq:ch01:spectral}
A=\sum_{j=1}^m \lambda_j\,\Pi_j,\qquad \Pi_j\Pi_k=0\ (j\neq k),\qquad \sum_{j=1}^m\Pi_j=I,
\end{equation}
where $\Pi_j$ is the orthogonal projector onto the eigenspace of $\lambda_j$.
\end{theorem}

The decomposition makes any function of a symmetric matrix a function of
numbers. For a function $f$ defined on the eigenvalues,
\begin{equation}\label{eq:ch01:funcalc}
f(A)=\sum_{j=1}^m f(\lambda_j)\,\Pi_j .
\end{equation}
With $f(t)=t^{-1}$ this is the inverse, which needs every $\lambda_j\neq0$.
With $f(t)=\log t$ it is the matrix logarithm of \cref{sec:ch01:logdet}, which
needs every $\lambda_j>0$. With $f(t)=\sqrt t$ it is the square root, which
needs every $\lambda_j\ge0$.
Among positive semidefinite matrices the square root is unique. Only one PSD
matrix $B$ has $B^2=A$, and it is the one \cref{eq:ch01:funcalc} gives
\cite{hornjohnson2012}. Other square roots exist if PSD is not required, for
example $-A^{1/2}$, and non-symmetric factors such as the Cholesky factor $L$
with $LL\T=A$ are also used. They are not ``the'' square root.

Geometrically, a symmetric matrix stretches space along its eigenvectors by
its eigenvalues. \Cref{fig:ch01:circle} shows the matrix of
Example~\ref{ex:ch01:spectral} sending the unit circle to an ellipse whose axes are
the eigenvectors, with half-lengths equal to the eigenvalues.

\begin{figure}[tbp]
\centering
\begin{tikzpicture}[scale=0.5]
\draw[guide,->] (-6,0)--(6.4,0) node[lbl,right] {$x_1$};
\draw[guide,->] (0,-3.2)--(0,3.6) node[lbl,above] {$x_2$};
\draw[thick,pgray] (0,0) circle (1);
\draw[thick,pblue,rotate=26.565] (0,0) ellipse (6.000 and 1.000);
\draw[vec,porange,rotate=26.565] (0,0)--(6,0) node[lbl,above left] {$Av_1=6v_1$};
\draw[vec,porange,rotate=26.565] (0,0)--(0,1);
\node[lbl,porange] at (-1.05,1.35) {$Av_2=v_2$};
\node[lbl,pgray] at (0.9,-1.2) {unit circle};
\end{tikzpicture}
\caption{The matrix with rows $(5,2)$ and $(2,2)$ maps the unit circle to the
blue ellipse. The eigenvector directions are the only ones that are stretched
without turning, by $6$ along $(2,1)\T$ and by $1$ along $(1,-2)\T$.}
\label{fig:ch01:circle}
\end{figure}

Positive semidefiniteness, written $A\psd 0$, means $x\T Ax\ge 0$ for all $x$
(Primer~L, section~L.6). By \cref{eq:ch01:spectral} that holds exactly when
every eigenvalue is nonnegative, and the eigenvalues bound the form,
\begin{equation}\label{eq:ch01:rayleighbounds}
\lambda_{\min}(A)\,\norm{x}^2\ \le\ x\T Ax\ \le\ \lambda_{\max}(A)\,\norm{x}^2 .
\end{equation}
The level set $x\T Ax=1$ shows the difference between definite and
semidefinite at a glance (\cref{fig:ch01:levels}). For $A\pd0$ it is an ellipse
with half-axes $1/\sqrt{\lambda_i}$. When an eigenvalue is zero, the form does
not grow along that eigenvector at all. In the plane, with the other eigenvalue
positive, the ellipse opens into a pair of parallel lines, and in higher
dimension it opens into a cylinder. A read operator with a kernel has level sets of this second
kind. Moving along the kernel never changes what the consumer sees.

\begin{figure}[tbp]
\centering
\begin{tikzpicture}[scale=1.1]
\begin{scope}
\clip (-2,-2) rectangle (2,2);
\draw[guide,->] (-2,0)--(2,0);
\draw[guide,->] (0,-2)--(0,2);
\draw[thick,porange,rotate=26.565] (0,0) ellipse (0.408 and 1);
\end{scope}
\node[lbl] at (0,-2.35) {$A\pd0$, rows $(5,2)$, $(2,2)$};
\begin{scope}[xshift=5cm]
\clip (-2,-2) rectangle (2,2);
\draw[guide,->] (-2,0)--(2,0);
\draw[guide,->] (0,-2)--(0,2);
\draw[thick,porange] (-1,2)--(3,-2);
\draw[thick,porange] (-3,2)--(1,-2);
\draw[vec,pred] (0,0)--(1.2,-1.2) node[lbl,right] {$\ker A$};
\end{scope}
\node[lbl] at (5,-2.35) {$A=gg\T\psd0$, $g=(1,1)\T$};
\end{tikzpicture}
\caption{The level set $x\T Ax=1$ is an ellipse when $A$ is positive definite
and a pair of parallel lines $x_1+x_2=\pm1$ when $A$ is semidefinite with a
kernel. Along the kernel direction the form stays constant forever.}
\label{fig:ch01:levels}
\end{figure}

The ratio of the two bounds in \cref{eq:ch01:rayleighbounds} is the
\emph{condition number} $\kappa(A)=\lambda_{\max}/\lambda_{\min}$ of a positive
definite matrix. For a covariance it says how elongated the cloud is. A
covariance with $\kappa=100$ has one direction with a hundred times the
variance of another, so its whitening $\Sigma^{-1/2}$ stretches the thin
direction by a factor of ten relative to the thick one. Any error made before
whitening, for example by estimating $\Sigma$ from too few samples, is
amplified along that thin direction. A large condition number is a warning
that a whitened quantity rests on the smallest eigenvalues, which are the
least reliably estimated. How far eigenvalues and eigenvectors move when a
matrix is perturbed is the subject of \cref{ch13}.

Comparing two symmetric matrices uses the same test on their difference.

\begin{definition}[Loewner order]\label{def:ch01:loewner}
For symmetric $A,B$ of the same size,
\begin{equation}\label{eq:ch01:loewner}
A\psd B \iff A-B\psd 0 \iff x\T Ax\ge x\T Bx\ \text{for every } x .
\end{equation}
\end{definition}

The relation is reflexive, antisymmetric and transitive, so it is an order. It
is only a \emph{partial} order. Two matrices can be incomparable, each larger
than the other along some direction. An error covariance can be smaller than
another along the directions one consumer reads and larger along the
directions a second consumer reads. That is why the programme compares errors
through a scalar such as $\tr(\PC\Sd)$, and why a statement of the form
$\Sigma_1\psd\Sigma_2$ is a much stronger claim, a win for every consumer at
once.

For covariances there is a picture. Draw each $\Sigma$ as its ellipse
$\{\Sigma^{1/2}u:\norm{u}\le1\}$, whose half-axes are $\sqrt{\lambda_i}$ along the
eigenvectors. Then $\Sigma_1\psd\Sigma_2$ exactly when the ellipse of
$\Sigma_2$ fits inside the ellipse of $\Sigma_1$. Two ellipses that cross
belong to incomparable matrices (\cref{fig:ch01:loewner}).

\begin{figure}[tbp]
\centering
\begin{tikzpicture}[scale=1.15]
\draw[guide] (-2.1,0)--(2.1,0); \draw[guide] (0,-2.1)--(0,2.1);
\draw[thick,pblue] (0,0) circle (1);
\draw[thick,porange,rotate=31.717] (0,0) ellipse (1.902 and 1.176);
\node[lbl,pblue] at (0.35,-0.7) {$I$};
\node[lbl,porange] at (1.55,1.55) {$A$};
\node[lbl] at (0,-2.45) {$A\psd I$, nested};
\begin{scope}[xshift=5cm]
\draw[guide] (-2.1,0)--(2.1,0); \draw[guide] (0,-2.1)--(0,2.1);
\draw[thick,pblue] (0,0) ellipse (1.732 and 1.000);
\draw[thick,porange,rotate=45] (0,0) ellipse (1.732 and 1.000);
\fill[pred] (1.732,0) circle (1.6pt);
\fill[pred] (1.225,1.225) circle (1.6pt);
\node[lbl,pblue] at (1.85,-0.65) {$\diag(3,1)$};
\node[lbl,porange] at (-1.75,1.45) {$\begin{pmatrix}2&1\\1&2\end{pmatrix}$};
\node[lbl] at (0,-2.45) {incomparable, crossing};
\end{scope}
\end{tikzpicture}
\caption{Covariance ellipses make the Loewner order visible. On the left the
ellipse of $A$, with rows $(3,1)$ and $(1,2)$, contains the unit circle, so
$A\psd I$. On the right each ellipse pokes out of the other, so neither matrix
dominates. The red points mark where each reaches farther than the other.}
\label{fig:ch01:loewner}
\end{figure}

Three facts carry most of the weight in later chapters.
\begin{align}
A\psd B &\implies C\T AC\psd C\T BC\ \text{for any } C, \label{eq:ch01:congruence}\\
A\psd B,\ P\psd 0 &\implies \tr(PA)\ge\tr(PB), \label{eq:ch01:tracemono}\\
A\psd B\pd 0 &\implies B^{-1}\psd A^{-1}. \label{eq:ch01:inversemono}
\end{align}
The first is $x\T C\T(A-B)Cx=(Cx)\T(A-B)(Cx)\ge0$. The second follows from
\cref{sec:ch01:trace}, since $\tr(P(A-B))\ge0$ when both factors are PSD. The
third has a short proof that uses whitening, given below. The second has a
converse. If $\tr(PA)\ge\tr(PB)$ for every $P\psd0$, then choosing $P=xx\T$
gives $x\T Ax\ge x\T Bx$ for every $x$. So
\begin{equation}\label{eq:ch01:everyreader}
A\psd B\iff \tr(PA)\ge\tr(PB)\ \text{for every } P\psd0 .
\end{equation}
In the language of the programme, one error covariance is below another in
the Loewner order exactly when every consumer, whatever its read operator,
finds it no worse. The ``if'' direction needs every rank-one reader $xx\T$.
Agreement among a few given consumers does not give the Loewner order. The
identity and the matrix with rows $(1,0.9)$ and $(0.9,1)$ tie for both
coordinate readers $e_1e_1\T$ and $e_2e_2\T$, yet their difference has
eigenvalues $\pm0.9$, so they are incomparable.
The proof of the inverse rule goes through whitening. By \cref{eq:ch01:congruence} with
$C=B^{-1/2}$, the hypothesis gives $B^{-1/2}AB^{-1/2}\psd I$, so every
eigenvalue of $B^{-1/2}AB^{-1/2}$ is at least $1$. Its inverse
$B^{1/2}A^{-1}B^{1/2}$ then has every eigenvalue at most $1$, so
$B^{1/2}A^{-1}B^{1/2}\preceq I$, and applying \cref{eq:ch01:congruence} with
$C=B^{-1/2}$ once more gives $A^{-1}\preceq B^{-1}$.

Not every operation that is monotone on numbers is monotone in this order.
Squaring is the standard warning. From $A\psd B\psd 0$ it does not follow that
$A^2\psd B^2$.

\begin{example}[A square root, an incomparable pair and a failure of squaring]\label{ex:ch01:spectral}
Let $A=\begin{pmatrix}5&2\\2&2\end{pmatrix}$. Its trace is $7$ and its
determinant is $6$, so its eigenvalues are $6$ and $1$, with unit
eigenvectors $(2,1)\T/\sqrt5$ and $(1,-2)\T/\sqrt5$. The two eigenprojectors are
\[
\Pi_1=\tfrac15\begin{pmatrix}4&2\\2&1\end{pmatrix},\qquad
\Pi_2=\tfrac15\begin{pmatrix}1&-2\\-2&4\end{pmatrix},
\]
and $6\Pi_1+\Pi_2=A$. The square root is $\sqrt6\,\Pi_1+\Pi_2$, which has
diagonal entries $(4\sqrt6+1)/5\approx2.1596$ and $(\sqrt6+4)/5\approx1.2899$
and off-diagonal entry $(2\sqrt6-2)/5\approx0.5798$. Its square has top-left
entry $2.1596^2+0.5798^2\approx5.000$.

The matrices $\diag(2,0)$ and $I$ are incomparable. Their difference
$\diag(1,-1)$ has one positive and one negative eigenvalue, so
$\diag(2,0)$ is larger along $e_1$ and smaller along $e_2$.

For squaring, take $A=\begin{pmatrix}2&1\\1&1\end{pmatrix}$ and
$B=\begin{pmatrix}1&0\\0&0\end{pmatrix}$. Then
$A-B=\begin{pmatrix}1&1\\1&1\end{pmatrix}$ has eigenvalues $2$ and $0$, so
$A\psd B\psd 0$. But $A^2-B^2=\begin{pmatrix}4&3\\3&2\end{pmatrix}$ has
determinant $-1$, so it has a negative eigenvalue and $A^2\not\psd B^2$.

For inversion, take $A=\begin{pmatrix}3&1\\1&2\end{pmatrix}$ and $B=I$. Then
$A-I$ has trace $3$ and determinant $1$, so $A\psd I$. And
$A^{-1}=\tfrac15\begin{pmatrix}2&-1\\-1&3\end{pmatrix}$, so
$I-A^{-1}=\tfrac15\begin{pmatrix}3&1\\1&2\end{pmatrix}$, with determinant
$\tfrac{5}{25}=0.2>0$ and positive trace. So $I\psd A^{-1}$, as
\cref{eq:ch01:inversemono} predicts.
\end{example}

\inproject{\nolinkurl{geometric-observation/paper/observation-theory.tex}, appendix
``The price of a misidentified observer'', theorem ``Commission tax'', assumes
$m\hat P\preceq P\preceq M\hat P$ and its proof uses \cref{eq:ch01:tracemono}
in the form $P\preceq M\hat P\Rightarrow\tr(P\Sigma_\delta)\le M\tr(\hat
P\Sigma_\delta)$. In \nolinkurl{geometric-observation/paper/tit-rate-leakage/tit-rate-leakage.tex},
section ``Convex Structure'', the step $\Sigma\preceq\Sigma_T$ if and only if
$\Sigma^{-1}\succeq\Sigma_T^{-1}$ is \cref{eq:ch01:inversemono}.}

\buildson{Primer~L, sections~L.6 and L.7.}
\usedin{Every later section of this chapter. \Cref{ch02}, where the
Cram\'er--Rao bound is a Loewner inequality. \Cref{ch04}, semidefinite
constraints. \Cref{ch06}, the commission and omission theorems.}

% ---------------------------------------------------------------------------
\ompnexthead{Trace as weighted variance}
\section{The trace of a product is a weighted sum of variances}\label{sec:ch01:trace}

Chapter~0 of \emph{Data Mining as Observation} defines the observer distortion as
a trace. This section collects the identities that make the trace usable and
shows why the trace of a product, and not a product of traces, is the right
scalar.

The trace is linear, and it does not see the order of two factors,
\begin{equation}\label{eq:ch01:cyclic}
\tr(AB)=\sum_{i,j}A_{ij}B_{ji}=\tr(BA),\qquad \tr(ABC)=\tr(CAB)=\tr(BCA).
\end{equation}
The rule allows cyclic shifts only. In general $\tr(ABC)\neq\tr(ACB)$
(\cref{exr:ch01:cyclic}). Here $A$ and $B$ need not be square, only
conformable, and that case gives the identity behind every quadratic form,
\begin{equation}\label{eq:ch01:qftrace}
x\T Ax=\tr(x\T Ax)=\tr(A\,xx\T),\qquad \norm{x}^2=\tr(xx\T).
\end{equation}

Taking expectations of \cref{eq:ch01:qftrace} turns a random quadratic form
into a deterministic trace. Let $\delta$ be a random error with mean $\mu$,
covariance $\Sd$ and second moment $M_\delta=\E[\delta\delta\T]=\Sd+\mu\mu\T$.
Linearity of trace and expectation gives, for a fixed matrix $P$,
\begin{equation}\label{eq:ch01:expqf}
\E\big[\delta\T P\delta\big]=\tr\big(P\,M_\delta\big)=\tr\big(P\,\Sd\big)+\mu\T P\mu .
\end{equation}
This is equation~0.10 of \emph{Data Mining as Observation}, with the mean term written out.
A biased error pays $\mu\T P\mu$ on top of its covariance term.

\begin{definition}[Observer distortion]\label{def:ch01:dO}
For a read operator $\PC\psd 0$ and a centred error with covariance $\Sd$,
\begin{equation}\label{eq:ch01:dO}
\dO=\tr(\PC\Sd)=\E\big[\delta\T\PC\,\delta\big].
\end{equation}
\end{definition}

The definition is a weighted mean squared error. Weighting squared error by a
PSD matrix is standard in estimation and in rate--distortion theory. The
programme takes the weight from the consumer's Jacobian, as that book does, and
\cref{ch06} gives the history of both the weighted error and that choice of
weight. The trace has a reading that explains the name.
Write $\PC=\sum_i p_i v_iv_i\T$ by the spectral theorem. Then
\begin{equation}\label{eq:ch01:weightedvar}
\tr(\PC\Sd)=\sum_i p_i\,v_i\T\Sd v_i=\sum_i p_i\Var\big(v_i\T\delta\big).
\end{equation}
The distortion is a sum of the error's variances along the consumer's
eigendirections, each weighted by how strongly the consumer reads that
direction. Setting $\PC=I$ recovers $\tr\Sd$, the total variance, which is
the identity reader. A direction in $\ker\PC$ contributes nothing however
large the error is there.

Two PSD matrices always have a nonnegative trace product, and it vanishes
only in a specific way,
\begin{equation}\label{eq:ch01:tracepos}
\tr(P\Sigma)=\tr\big(P^{1/2}\Sigma P^{1/2}\big)=\big\|\Sigma^{1/2}P^{1/2}\big\|_F^2\ \ge 0,
\qquad \tr(P\Sigma)=0 \iff P\Sigma=0 .
\end{equation}
The equality case says the error lives entirely in the kernel of the reader.

\begin{example}[Same total error, different distortion]\label{ex:ch01:trace}
Let the consumer have $\PC=\begin{pmatrix}2&1\\1&2\end{pmatrix}$, with
eigenvalue $3$ along $(1,1)\T/\sqrt2$ and $1$ along $(1,-1)\T/\sqrt2$. Compare
two centred errors with the same total variance $4$,
\[
\Sigma_a=\begin{pmatrix}1&0\\0&3\end{pmatrix},\qquad
\Sigma_b=\begin{pmatrix}2&1\\1&2\end{pmatrix}.
\]
For $\Sigma_a$, the variance along $(1,1)\T/\sqrt2$ is $(1+3)/2=2$ and along
$(1,-1)\T/\sqrt2$ is also $2$, so \cref{eq:ch01:weightedvar} gives
$3\cdot2+1\cdot2=8$. Directly, $\PC\Sigma_a=\begin{pmatrix}2&3\\1&6\end{pmatrix}$
has trace $8$. For $\Sigma_b$, the variances along the two eigendirections are
$3$ and $1$, so the distortion is $3\cdot3+1\cdot1=10$, and
$\tr(\PC\Sigma_b)=\tr\begin{pmatrix}5&4\\4&5\end{pmatrix}=10$. The identity
reader scores the two errors equal at $4$. This consumer scores $\Sigma_b$
worse, because $\Sigma_b$ puts its error along the direction the consumer reads
most. Adding a bias $\mu=(1,0)\T$ to $\Sigma_a$ raises the distortion by
$\mu\T\PC\mu=2$, to $10$.
\end{example}

\inproject{\nolinkurl{geometric-observation/paper/observation-theory.tex}, section GO-2
on distortion (label \texttt{sec:distortion}), compares codes at matched
reconstruction error and reports that $\tr(\PC\Sd)$ tracks which code the
consumer prefers while reconstruction error does not, with the claim status
\texttt{[demonstrated]}. Example~\ref{ex:ch01:trace} is the two-dimensional version of
that dissociation.}

\buildson{Chapter~0 of \emph{Data Mining as Observation}, sections~0.3 and 0.5,
equations~0.3, 0.4 and 0.10. Primer~S, section~S.2 (linearity of expectation).}
\usedin{\Cref{ch02}, where $\tr(\PC\Sigma)$ is evaluated on Kalman covariances.
\Cref{ch03}, consumer-relative rate--distortion. \Cref{ch06} throughout. \Cref{ch07},
where a quantizer's error covariance is scored by $\dO$.}

% ---------------------------------------------------------------------------
\ompnexthead{Whitening}
\section{Whitening makes the source isotropic, and the whitened operator carries the reader into those coordinates}\label{sec:ch01:whiten}

Primer~L, section~L.8, shows that the singular value decomposition of a
centred data matrix $X\in\R^{n\times d}$ gives the principal components,
\begin{equation}\label{eq:ch01:svd}
X=U\Sigma_{\!s}V\T,\qquad \Sigma=\tfrac1n X\T X=V\Lambda V\T,\qquad \Lambda=\tfrac1n\Sigma_{\!s}^2 .
\end{equation}
Principal component analysis keeps the top $k$ columns $V_k$ of $V$. The
projector $V_kV_k\T$ keeps the variance $\tr(V_k\T\Sigma V_k)=\lambda_1+\dots+\lambda_k$,
and \cref{sec:ch01:pencil} shows that no other $k$-dimensional subspace keeps more.
That is the identity reader's answer. A consumer that reads other directions
wants a different subspace, and the tool that compares a reader with a source
is whitening.

\begin{definition}[Whitening]\label{def:ch01:whiten}
For $\Sigma\pd 0$, a matrix $W$ \emph{whitens} $\Sigma$ when $W\Sigma W\T=I$.
The symmetric choice and the general solution are
\begin{equation}\label{eq:ch01:whiten}
W=\Sigma^{-1/2},\qquad\text{and every whitening matrix is } W=Q\,\Sigma^{-1/2}\ \text{with } Q\T Q=I .
\end{equation}
\end{definition}

To see the second claim, note that $Q=W\Sigma^{1/2}$ satisfies
$QQ\T=W\Sigma W\T=I$, so it is orthogonal. Whitening is therefore unique only
up to a rotation, a point \cref{sec:ch01:orth} returns to. The choice
$Q=V\T$ gives PCA whitening $\Lambda^{-1/2}V\T$, which rotates into the
principal axes and then rescales. The choice $Q=I$ gives the symmetric whitening
$\Sigma^{-1/2}$, which rescales without leaving the original orientation.

Now let a source $x$ have covariance $\Sigma_x\pd0$ and write it as
$x=\mu+\Sigma_x^{1/2}z$, so that $z$ has covariance $I$. An error $\delta_z$ in
whitened coordinates becomes $\delta_x=\Sigma_x^{1/2}\delta_z$ in source
coordinates. A consumer with read operator $P$ in source coordinates sees
\begin{equation}\label{eq:ch01:whitenedop}
\delta_x\T P\,\delta_x=\delta_z\T\tilde P\,\delta_z,\qquad \tilde P=\Sigma_x^{1/2}P\,\Sigma_x^{1/2}.
\end{equation}
The matrix $\tilde P$ is the \emph{whitened operator}. Its eigenvalues weigh
how much a consumer cares about each direction against how much variance the
source has there, and they are what reverse water-filling in \cref{ch03} acts on.
Three facts make it easy to use,
\begin{equation}\label{eq:ch01:whitenfacts}
\tr(P\Sigma_x)=\tr\tilde P,\qquad
\operatorname{spec}\tilde P=\operatorname{spec}(P\Sigma_x),\qquad
\ker\tilde P=\Sigma_x^{-1/2}\ker P .
\end{equation}
The first is the cyclic rule. The second holds because
$P\Sigma_x=\Sigma_x^{-1/2}\tilde P\,\Sigma_x^{1/2}$ is similar to $\tilde P$,
even though $P\Sigma_x$ itself is not symmetric. The same argument shows that
$P^{1/2}\Sigma_xP^{1/2}$ has the same eigenvalues as well, so papers that use
either symmetric form are diagonalizing the same spectrum. The third holds
because $\tilde Pv=0$ exactly when $P\Sigma_x^{1/2}v=0$. The whitened operator
is \emph{not} $P$ with its eigenvalues rescaled. Its eigenvectors differ from
those of $P$ and of $\Sigma_x$ whenever the two do not commute.

\begin{example}[A reader of the sum, whitened]\label{ex:ch01:whiten}
Let $\Sigma_x=\diag(4,1)$ and let the consumer read $x_1+x_2$, so $g=(1,1)\T$
and $P=gg\T$, with every entry $1$. The eigenvalues of $P$ are $2$ and $0$. The
whitened operator is
\[
\tilde P=\diag(2,1)\begin{pmatrix}1&1\\1&1\end{pmatrix}\diag(2,1)
=\begin{pmatrix}4&2\\2&1\end{pmatrix}=(2,1)\T(2,1),
\]
with eigenvalue $5$ along $(2,1)\T/\sqrt5$ and $0$ along $(1,-2)\T/\sqrt5$.
The trace $\tr(P\Sigma_x)=4+1=5$ matches. The nonsymmetric product
$P\Sigma_x=\begin{pmatrix}4&1\\4&1\end{pmatrix}$ also has trace $5$ and
determinant $0$, so the same eigenvalues. The whitened kernel
$(1,-2)\T$ maps back to $\Sigma_x^{1/2}(1,-2)\T=(2,-2)\T$, which spans
$\ker P$, as the third identity in \cref{eq:ch01:whitenfacts} says. In whitened
coordinates the consumer's top direction leans toward the first coordinate,
because that is where the source has more variance to spoil.
\end{example}

\Cref{fig:ch01:whiten} shows whitening on a sample. Sixty draws from
$\Normal(0,\Sigma)$, with $\Sigma$ the matrix of Example~\ref{ex:ch01:spectral}, form
a tilted cloud. Multiplying each point by $\Sigma^{-1/2}$ turns the cloud
round, and the covariance ellipse becomes the unit circle.

\begin{figure}[tbp]
\centering
\begin{tikzpicture}
\begin{axis}[primer, width=0.47\linewidth, height=0.42\linewidth, axis equal image,
  xmin=-5.5, xmax=5.5, ymin=-4.2, ymax=4.2, title={\small before, $x$}]
\addplot[only marks, mark=*, mark size=1.1pt, pblue] coordinates {(2.93,-2.11) (0.57,-0.49) (-1.10,-0.54) (-4.50,-1.47) (0.06,3.78) (0.28,-0.32) (-0.99,-1.02) (-2.51,-1.12) (0.90,-0.03) (1.95,0.30) (0.95,2.01) (0.88,-0.34) (-0.08,0.59) (4.02,0.77) (0.06,1.15) (-2.08,-0.89) (2.24,1.26) (0.59,0.92) (-5.52,-0.32) (-3.04,-2.71) (1.00,1.06) (-1.58,-1.65) (0.03,-0.05) (3.47,1.78) (1.06,1.55) (-0.98,-1.31) (1.60,1.09) (-0.92,-1.13) (-0.95,-3.08) (1.78,1.03) (-3.50,-0.87) (-1.64,0.42) (-4.92,-2.36) (2.20,1.90) (-4.95,-1.89) (0.36,-0.60) (2.74,-0.61) (0.16,-1.15) (3.02,0.79) (-1.80,-2.43) (4.07,1.95) (2.29,1.90) (0.38,-0.62) (-2.19,-1.50) (2.11,-1.09) (-1.47,-0.76) (0.82,0.87) (-3.70,-2.96) (0.69,1.56) (1.76,0.72) (-0.51,0.19) (-0.05,0.91) (-1.64,-0.29) (0.08,0.64) (1.96,3.42) (4.10,2.80) (-4.60,-1.62) (-1.01,0.33) (-4.24,0.19) (1.55,-1.06)};
\addplot[thick,porange,domain=0:360,samples=91] ({2.449*cos(x)*cos(26.565)-sin(x)*sin(26.565)},{2.449*cos(x)*sin(26.565)+sin(x)*cos(26.565)});
\end{axis}
\end{tikzpicture}\hfill
\begin{tikzpicture}
\begin{axis}[primer, width=0.47\linewidth, height=0.42\linewidth, axis equal image,
  xmin=-5.5, xmax=5.5, ymin=-4.2, ymax=4.2, title={\small after, $\Sigma^{-1/2}x$}]
\addplot[only marks, mark=*, mark size=1.1pt, pgreen] coordinates {(2.04,-2.56) (0.42,-0.57) (-0.45,-0.22) (-2.02,-0.23) (-0.87,3.32) (0.23,-0.35) (-0.28,-0.67) (-1.06,-0.39) (0.48,-0.24) (0.96,-0.20) (0.02,1.55) (0.55,-0.51) (-0.18,0.54) (1.94,-0.27) (-0.24,1.00) (-0.89,-0.29) (0.88,0.58) (0.09,0.67) (-2.83,1.02) (-0.96,-1.67) (0.28,0.70) (-0.44,-1.08) (0.03,-0.05) (1.41,0.75) (0.19,1.11) (-0.21,-0.93) (0.58,0.58) (-0.21,-0.78) (0.23,-2.49) (0.69,0.49) (-1.64,0.06) (-0.96,0.76) (-2.03,-0.91) (0.71,1.16) (-2.16,-0.50) (0.33,-0.61) (1.59,-1.19) (0.35,-1.05) (1.41,-0.02) (-0.37,-1.72) (1.68,0.75) (0.75,1.14) (0.35,-0.64) (-0.80,-0.80) (1.37,-1.46) (-0.60,-0.32) (0.22,0.58) (-1.25,-1.73) (-0.00,1.21) (0.76,0.22) (-0.32,0.29) (-0.24,0.82) (-0.79,0.13) (-0.11,0.54) (0.22,2.55) (1.50,1.50) (-2.04,-0.34) (-0.61,0.53) (-2.28,1.17) (1.07,-1.30)};
\addplot[thick,porange,domain=0:360,samples=91] ({cos(x)},{sin(x)});
\end{axis}
\end{tikzpicture}
\caption{Whitening turns a correlated cloud into a round one. The orange curve is
the covariance ellipse, half-axes $\sqrt6$ and $1$ before and the unit circle
after. Distances in the right panel are Mahalanobis distances in the left.}
\label{fig:ch01:whiten}
\end{figure}

\inproject{\nolinkurl{geometric-observation/paper/observation-theory.tex} defines
$\tilde P:=\Sigma_x^{1/2}P\Sigma_x^{1/2}$ as ``the whitened operator'' in appendix
``The price of a misidentified observer''. The theorem ``Omission floor'' projects
onto the \emph{whitened} kernel $\Sigma_x^{-1/2}\ker\hat P$, and its proof explains
why the naive kernel $\ker\hat P$ is the wrong one when $\Sigma_x\neq I$. The
Berger-reduction proposition (label \texttt{prop:rd}) in section ``A
consumer-relative rate--distortion function'' diagonalizes
$\PC^{1/2}\Sigma_x\PC^{1/2}$, which has the same spectrum.}

\buildson{Primer~L, sections~L.7 and L.8. Chapter~0 of \emph{Data Mining as
Observation}, equation~0.6.}
\usedin{\Cref{ch02} (whitening a Gaussian), \cref{ch03} (reverse water-filling on
the eigenvalues of $\tilde P$), \cref{ch06} (omission floor).}

% ---------------------------------------------------------------------------
\ompnexthead{Pencils and Rayleigh quotients}
\section{Two quadratic forms can be diagonalized together, and the Rayleigh quotient ranks the directions}\label{sec:ch01:pencil}

Many questions in the programme ask for the direction that does best on one
quadratic form relative to another. Which one-dimensional description of a
source buys the most reduction in a consumer's error per unit of leaked
information? Which direction carries the most consumer-weighted variance per
unit of source variance? When the form in the denominator is positive
definite, the answer is a generalized eigenvector. When it is only
semidefinite, the ratio can be unbounded along its kernel.

\begin{definition}[Matrix pencil and generalized eigenproblem]\label{def:ch01:pencil}
For symmetric $A$ and symmetric $B\pd 0$, the \emph{pencil} $(A,B)$ has
generalized eigenpairs $(\mu,v)$ with
\begin{equation}\label{eq:ch01:geneig}
Av=\mu\,Bv,\qquad \det(A-\mu B)=0 .
\end{equation}
\end{definition}

Substituting $v=B^{-1/2}w$ turns the pencil into an ordinary symmetric
eigenproblem for the matrix $B^{-1/2}AB^{-1/2}$. So the $\mu$ are real, and the
eigenvectors can be chosen orthonormal \emph{in the $B$-inner product}.
Collecting them in a matrix $V$ diagonalizes both forms at once,
\begin{equation}\label{eq:ch01:simdiag}
V\T BV=I,\qquad V\T AV=\diag(\mu_1,\dots,\mu_d).
\end{equation}
This is simultaneous diagonalization. In the coordinates $c=V^{-1}x$ both forms
are sums of squares, $x\T Bx=\sum_ic_i^2$ and $x\T Ax=\sum_i\mu_ic_i^2$. This
construction needs $B\pd0$. That is sufficient, not necessary. Two PSD
matrices $A$ and $B$ with $A+B\pd0$ can still be diagonalized together, by
applying \cref{eq:ch01:simdiag} to the pencil $(A,A+B)$. Then $V\T(A+B)V=I$
and $V\T AV$ is diagonal, so $V\T BV=I-V\T AV$ is diagonal too. Without any
definiteness it can fail. For $A=\diag(1,-1)$ and $B$ with rows $(0,1)$ and
$(1,0)$, $\det(A-\mu B)=-1-\mu^2$ has no real root, so no real congruence
diagonalizes both.

The ratio of the two forms is the \emph{Rayleigh quotient} of the pencil,
\begin{equation}\label{eq:ch01:rayleigh}
\rho(x)=\frac{x\T Ax}{x\T Bx}=\frac{\sum_i\mu_ic_i^2}{\sum_ic_i^2},\qquad
\mu_{\min}\le\rho(x)\le\mu_{\max}.
\end{equation}
The second form of $\rho$ shows it is a weighted average of the $\mu_i$, so it
lies between the extreme generalized eigenvalues and reaches them at the
corresponding eigenvectors. Its gradient is
$\nabla\rho=\tfrac{2}{x\T Bx}\big(Ax-\rho(x)Bx\big)$, so the stationary points
of $\rho$ are exactly the generalized eigenvectors. The intermediate
eigenvalues have a variational description too.

\begin{theorem}[Courant--Fischer]\label{thm:ch01:cf}
With $\mu_1\ge\dots\ge\mu_d$ the generalized eigenvalues of $(A,B)$,
\begin{equation}\label{eq:ch01:courant}
\mu_k=\max_{\dim\mathcal S=k}\ \min_{0\neq x\in\mathcal S}\rho(x)
=\min_{\dim\mathcal S=d-k+1}\ \max_{0\neq x\in\mathcal S}\rho(x).
\end{equation}
\end{theorem}

With $B=I$ this is the classical statement for eigenvalues of a symmetric
matrix \cite{hornjohnson2012}. The pencil version follows by the substitution
$x=B^{-1/2}w$, which maps subspaces of dimension $k$ to subspaces of dimension
$k$. Asking for the best $k$-dimensional subspace as a whole, rather than the
$k$-th best direction, gives the trace form.

\begin{theorem}[Ky Fan's maximum principle]\label{thm:ch01:kyfan}
For symmetric $A$ with eigenvalues $\lambda_1\ge\dots\ge\lambda_d$,
\begin{equation}\label{eq:ch01:kyfan}
\max_{U\in\R^{d\times k},\ U\T U=I_k}\tr\big(U\T AU\big)=\lambda_1+\dots+\lambda_k ,
\end{equation}
attained when the columns of $U$ span the top-$k$ eigenspace. For the pencil
$(A,B)$ the same holds with the constraint $U\T BU=I_k$ and the generalized
eigenvalues $\mu_1+\dots+\mu_k$.
\end{theorem}

The result is due to Ky Fan \cite{fan1949}. With $A=\Sigma$ it is the PCA
statement of \cref{sec:ch01:whiten}. Note that $\tr(U\T AU)=\tr(A\,UU\T)$ is the
trace of $A$ against the projector $UU\T$, the form of \cref{sec:ch01:trace}.

\begin{example}[A pencil whose eigenvectors are not orthogonal]\label{ex:ch01:pencil}
Let $A=\begin{pmatrix}2&2\\2&8\end{pmatrix}$ and $B=\diag(1,4)$. Then
$\det(A-\mu B)=(2-\mu)(8-4\mu)-4=4\mu^2-16\mu+12$, with roots $\mu=3$ and
$\mu=1$. For $\mu=3$, $A-3B=\begin{pmatrix}-1&2\\2&-4\end{pmatrix}$ and
$v_1\propto(2,1)\T$. For $\mu=1$, $A-B=\begin{pmatrix}1&2\\2&4\end{pmatrix}$ and
$v_2\propto(2,-1)\T$. Both have $v\T Bv=4+4=8$, so the $B$-normalized vectors
are $(2,\pm1)\T/(2\sqrt2)$. They are $B$-orthogonal, since
$(2,1)B(2,-1)\T=4-4=0$, but not orthogonal in the dot product, since
$(2,1)\cdot(2,-1)=3$. With $V$ holding the normalized vectors, $V\T BV=I$ and
$V\T AV=\diag(3,1)$. The check on the first entry is $A(2,1)\T=(6,12)\T$,
then $(2,1)(6,12)\T/8=24/8=3$. At $x=e_1$ the Rayleigh quotient is
$2/1=2$, between $1$ and $3$. The ordinary eigenvalues of $A$ are
$5\pm\sqrt{13}$, about $8.606$ and $1.394$, which shows that the pencil, and not
$A$ alone, answers the relative question. \Cref{fig:ch01:rayleigh} traces the
Rayleigh quotient around the half circle of directions.

For Ky Fan, let $A=\begin{pmatrix}3&1&0\\1&3&0\\0&0&1\end{pmatrix}$, with
eigenvalues $4$, $2$ and $1$. The best two-dimensional trace is $4+2=6$. The
subspace spanned by $e_1,e_2$ is the top eigenspace and gives
$\tr(U\T AU)=3+3=6$. The subspace spanned by $e_1,e_3$ gives $3+1=4$.
\end{example}

\begin{figure}[tbp]
\centering
\begin{tikzpicture}
\begin{axis}[primer, xmin=0, xmax=180, ymin=0.6, ymax=3.4,
  xtick={0,45,90,135,180}, xlabel={\small direction angle of $x$ (degrees)},
  ylabel={\small $\rho(x)$}]
\addplot[thick,pblue] coordinates {(0,2.0000) (5,2.3396) (10,2.6273) (15,2.8327) (20,2.9516) (25,2.9976) (30,2.9897) (35,2.9459) (40,2.8795) (45,2.8000) (50,2.7135) (55,2.6238) (60,2.5329) (65,2.4423) (70,2.3523) (75,2.2632) (80,2.1750) (85,2.0873) (90,2.0000) (95,1.9127) (100,1.8250) (105,1.7368) (110,1.6477) (115,1.5577) (120,1.4671) (125,1.3762) (130,1.2865) (135,1.2000) (140,1.1205) (145,1.0541) (150,1.0103) (155,1.0024) (160,1.0484) (165,1.1673) (170,1.3727) (175,1.6604) (180,2.0000)};
\addplot[guide] coordinates {(0,3) (180,3)};
\addplot[guide] coordinates {(0,1) (180,1)};
\addplot[only marks, mark=*, mark size=2pt, pgreen] coordinates {(26.565,3) (153.435,1)};
\node[lbl,pgreen,above right] at (axis cs:26.565,3) {$v_1\propto(2,1)$, $\mu_1=3$};
\node[lbl,pgreen,above left] at (axis cs:153.435,1) {$v_2\propto(2,-1)$, $\mu_2=1$};
\node[lbl,pblue,below right] at (axis cs:0,2) {$e_1$, $\rho=2$};
\end{axis}
\end{tikzpicture}
\caption{The Rayleigh quotient of the pencil with $A$ rows $(2,2)$, $(2,8)$ and
$B=\diag(1,4)$, over all directions in the plane. It peaks at the generalized
eigenvalue $3$ and bottoms out at $1$, and the two extreme directions are
$26.6^\circ$ and $153.4^\circ$, which are not at right angles.}
\label{fig:ch01:rayleigh}
\end{figure}

\inproject{In \nolinkurl{geometric-observation/paper/tit-rate-leakage/tit-rate-leakage.tex},
section ``The One-Dimensional Case'', the theorem labelled \texttt{thm:onedim}
gives the minimum leakage as $\tfrac12\log(1+\Delta/\mu_{\max})$, and its proof
reduces the minimization to maximizing ``the Rayleigh quotient
$b\T A(\Delta)b/b\T\Kc b$ of a symmetric pencil''. Section ``Directions and
Misalignment'' states that, with no second party, the best one-dimensional
descriptions read the top generalized eigenspace of the pencil
$(\Sigma_TQ\Sigma_T,\Sigma_T)$. The observation-theory paper's ``commuting
observers'' case in section ``Successive refinement across observers'' assumes
that $\Sigma_x,P_1,P_2$ are simultaneously diagonalizable.}

\buildson{\Cref{sec:ch01:spectral,sec:ch01:whiten}. Primer~L, section~L.7,
the PCA paragraph.}
\usedin{\Cref{ch03}, where water-filling acts on simultaneously diagonal forms.
\Cref{ch04}, Lagrange multipliers for the Rayleigh quotient. \Cref{ch06}, the pencil
characterization of leakage.}


% ---------------------------------------------------------------------------
\ompnexthead{Schur complements}
\section{The Schur complement is what remains of a block after the other block is eliminated}\label{sec:ch01:schur}

Gaussian elimination on a matrix split into blocks produces one matrix again
and again. It is the conditional covariance of \cref{ch02}, the error covariance
of a Kalman update and the certificate in a semidefinite constraint.

\begin{definition}[Schur complement]\label{def:ch01:schur}
For a symmetric block matrix with $A$ invertible,
\begin{equation}\label{eq:ch01:schur}
M=\begin{pmatrix}A&B\\B\T&D\end{pmatrix},\qquad M/A=S=D-B\T A^{-1}B .
\end{equation}
\end{definition}

Subtracting $B\T A^{-1}$ times the first block row from the second clears the
lower-left block, and doing the same with columns clears the upper-right. The
result is a block factorization,
\begin{equation}\label{eq:ch01:ldu}
M=\begin{pmatrix}I&0\\B\T A^{-1}&I\end{pmatrix}
\begin{pmatrix}A&0\\0&S\end{pmatrix}
\begin{pmatrix}I&A^{-1}B\\0&I\end{pmatrix}.
\end{equation}
The outer factors are triangular with unit diagonal, so they have determinant
one and are easy to invert. Three consequences follow at once,
\begin{gather}
\det M=\det A\cdot\det S, \label{eq:ch01:schurdet}\\
M^{-1}=\begin{pmatrix}A^{-1}+A^{-1}BS^{-1}B\T A^{-1}&-A^{-1}BS^{-1}\\-S^{-1}B\T A^{-1}&S^{-1}\end{pmatrix}, \label{eq:ch01:schurinv}
\end{gather}
and the positive definiteness test below. The inverse formula needs $S$ to be
invertible as well. The lower-right block of $M^{-1}$ is then
$S^{-1}$. That fact is the bridge to \cref{ch02}, where the inverse covariance
holds the conditional precisions.

\begin{theorem}[Schur complement test]\label{thm:ch01:schurpsd}
If $A\pd0$, then
\begin{equation}\label{eq:ch01:schurpsd}
M\psd 0 \iff S\psd0,\qquad M\pd0\iff S\pd0 .
\end{equation}
\end{theorem}

\begin{proof}
The outer factors in \cref{eq:ch01:ldu} are transposes of each other, so
$M=L\,\diag(A,S)\,L\T$ with $L$ invertible. By \cref{eq:ch01:congruence},
congruence by an invertible matrix preserves both $\psd0$ and $\pd0$ in both
directions. And $\diag(A,S)$ is PSD (PD) exactly when $S$ is, given $A\pd0$.
\end{proof}

Eliminating in the other order, with $D$ invertible, gives the other Schur
complement $A-BD^{-1}B\T$. Equating the upper-left blocks of the two expressions
for $M^{-1}$ gives the matrix inversion lemma, often called the Woodbury
identity. It holds when $A$, $D$ and both Schur complements are invertible,
\begin{equation}\label{eq:ch01:woodbury}
(A-BD^{-1}B\T)^{-1}=A^{-1}+A^{-1}B\big(D-B\T A^{-1}B\big)^{-1}B\T A^{-1}.
\end{equation}
\Cref{ch02} uses it to pass between the covariance and the information forms of
the Kalman update.

\begin{example}[A three-by-three matrix split one and two]\label{ex:ch01:schur}
Let
\[
M=\begin{pmatrix}4&2&2\\2&2&1\\2&1&3\end{pmatrix},\qquad A=4,\quad B=(2,2),\quad D=\begin{pmatrix}2&1\\1&3\end{pmatrix}.
\]
Then $B\T A^{-1}B=\tfrac14\begin{pmatrix}4&4\\4&4\end{pmatrix}$, every entry $1$,
and
\[
S=D-B\T A^{-1}B=\begin{pmatrix}1&0\\0&2\end{pmatrix}.
\]
So $\det M=4\cdot2=8$. Expanding along the first row gives the same number,
$4(6-1)-2(6-2)+2(2-4)=20-8-4=8$. Since $A=4>0$ and $S\pd0$, $M\pd0$. The
lower-right block of $M^{-1}$ is $S^{-1}=\diag(1,\tfrac12)$. Changing the
bottom-right entry of $M$ from $3$ to $\tfrac12$ changes $S$ to
$\diag(1,-\tfrac12)$, and the test says that $M$ is then not PSD.
\end{example}

\begin{figure}[tbp]
\centering
\begin{tikzpicture}[node distance=2.2cm]
\node[block] (M) {$M=\begin{pmatrix}A&B\\B\T&D\end{pmatrix}$};
\node[block, right=2.6cm of M] (Dg) {$\begin{pmatrix}A&0\\0&S\end{pmatrix}$};
\draw[flow] (M) -- node[lbl,above] {eliminate} node[lbl,below] {$L^{-1}(\cdot)L^{-\top}$} (Dg);
\node[block, fill=white, below=1.0cm of Dg, text width=5.2cm, align=left] (out)
  {$S=D-B\T A^{-1}B$\\[2pt]
   $\det M=\det A\cdot\det S$\\[2pt]
   $M\psd0\iff S\psd0$ \ (given $A\pd0$)\\[2pt]
   lower-right block of $M^{-1}$ is $S^{-1}$};
\draw[flow,pgreen] (Dg) -- (out);
\node[lbl,pgray,below=1.0cm of M, text width=4.2cm, align=center] (g)
  {in \cref{ch02}, $S$ is the covariance of the second block given the first};
\draw[flow,pgray] (out.west) -- (g.east);
\end{tikzpicture}
\caption{Block elimination turns $M$ into a block-diagonal matrix by a congruence
with unit-triangular factors. Everything about $M$ that the chapter needs can
then be read from $A$ and the Schur complement $S$.}
\label{fig:ch01:schur}
\end{figure}

\inproject{\nolinkurl{geometric-observation/paper/tit-rate-leakage/tit-rate-leakage.tex},
section ``Convex Structure'', turns a nonlinear matrix constraint into the linear
matrix inequality labelled \texttt{eq:lmi} ``by the Schur complement criterion''
and uses the Woodbury identity in the same paragraph. Appendix ``The Gaussian
Exhaustion Lemma'' computes a determinant by ``the Schur factorization in the
other order'', which is \cref{eq:ch01:schurdet}.}

\buildson{Primer~L, section~L.3 (determinants and inverses). \Cref{sec:ch01:spectral},
congruence.}
\usedin{\Cref{ch02}, Gaussian conditioning and the Kalman update. \Cref{ch04},
semidefinite programs and the max-det problem.}

% ---------------------------------------------------------------------------
\ompnexthead{Orthogonal matrices}
\section{Orthogonal matrices preserve every length, and a random one spreads a vector's energy across coordinates}\label{sec:ch01:orth}

Primer~L, section~L.5, defines an orthogonal matrix by $Q\T Q=I$. This section
collects what such matrices preserve, what they do not, and why vector
quantizers multiply by a random one before rounding.

\begin{definition}[Orthogonal group and rotations]\label{def:ch01:orth}
\begin{equation}\label{eq:ch01:orth}
O(d)=\{Q\in\R^{d\times d}: Q\T Q=I\},\qquad \det Q=\pm1,\qquad
SO(d)=\{Q\in O(d):\det Q=1\}.
\end{equation}
\end{definition}
The matrices in $SO(d)$ are the rotations. Those with determinant $-1$ include
a reflection. In the plane every rotation has the form
$R(\varphi)=\begin{pmatrix}\cos\varphi&-\sin\varphi\\ \sin\varphi&\cos\varphi\end{pmatrix}$.

An orthogonal change of coordinates preserves lengths, angles, traces and
eigenvalues, but it moves a covariance relative to any reader that stays put,
\begin{equation}\label{eq:ch01:orthinv}
\norm{Qx}=\norm{x},\qquad \tr(Q\Sigma Q\T)=\tr\Sigma,\qquad
\tr\big(P\,Q\Sigma Q\T\big)=\tr\big(Q\T PQ\,\Sigma\big)\neq\tr(P\Sigma)\ \text{in general}.
\end{equation}
The identity reader cannot tell a rotated error from the original. A consumer
with $P\neq cI$ can. One special case matters for quantizers. If an encoder
rotates $x$ to $Qx$, makes an error $\varepsilon$ in rotated coordinates with
covariance $\sigma^2I$, and rotates back, the error in source coordinates is
$Q\T\varepsilon$, whose covariance is $\sigma^2Q\T Q=\sigma^2I$ again. Isotropic
error stays isotropic, and its distortion for any reader is $\sigma^2\tr P$.

A \emph{Haar-distributed} random orthogonal matrix is one whose distribution is
unchanged by multiplying it by any fixed orthogonal matrix. The standard recipe
draws $G$ with independent $\Normal(0,1)$ entries, factors $G=QR$, and
corrects signs so that $R$ has a positive diagonal,
\begin{equation}\label{eq:ch01:haar}
G=QR,\qquad Q_{\mathrm{Haar}}=Q\,\diag\big(\operatorname{sign}r_{11},\dots,\operatorname{sign}r_{dd}\big).
\end{equation}
The correction matters. LAPACK-based QR routines such as \texttt{np.linalg.qr} do not guarantee a
positive diagonal of $R$, and without the correction the output is not exactly Haar
distributed \cite{mezzadri2007}. Invariance gives the property quantizers rely
on. For a fixed $x$, the matrix $C=\E[Qxx\T Q\T]$ satisfies $UCU\T=C$ for every
orthogonal $U$, and the only symmetric matrices with that property are
multiples of $I$. Its trace is $\norm{x}^2$, so
\begin{equation}\label{eq:ch01:spread}
\E\big[(Qx)(Qx)\T\big]=\frac{\norm{x}^2}{d}\,I .
\end{equation}
Averaged over the draw of $Q$, every coordinate of the rotated vector has
second moment $\norm{x}^2/d$, however the energy of $x$ was spread before. For
a single draw of $Q$ the squared coordinates scatter around this value, in the
way described next. A scalar quantizer designed for that
spread then fits every coordinate.

More is true for a unit vector $x$. The rotated vector $Qx$ is uniformly
distributed on the unit sphere, and in high dimension each of its
coordinates is close in distribution to $\Normal(0,1/d)$. That is why
turboquant-pro scales a single Lloyd--Max codebook for the standard normal by
$1/\sqrt d$ and uses it on every coordinate, with the vector's norm stored
separately. The rotation does not reduce any error by itself. It reshapes the
input so that one fixed scalar codebook suits all coordinates, and so that,
with high probability over the rotation, no coordinate carries an outlier the
codebook was not designed for. The price is
that coordinates of $Qx$ are no longer the original features, so a consumer
that reads a few original features now depends on all rotated coordinates at
once.

A dense random $Q$ costs $d^2$ multiplications to apply. The Hadamard transform
gives a structured alternative. Sylvester's construction is
\begin{equation}\label{eq:ch01:hadamard}
H_1=(1),\qquad H_{2m}=\begin{pmatrix}H_m&H_m\\H_m&-H_m\end{pmatrix},\qquad
H_dH_d\T=d\,I,\qquad R=\tfrac{1}{\sqrt d}\,H_d\,D_s ,
\end{equation}
where $d$ is a power of two and $D_s=\diag(s_1,\dots,s_d)$ holds independent
random signs. The matrix $R$ is orthogonal, and the recursion applies it in
$d\log_2d$ additions and subtractions. Every entry of $H_d$ is $\pm1$, so
$R e_j$ has every coordinate equal to $\pm1/\sqrt d$. Each coordinate of $Rx$
has second moment $\norm{x}^2/d$ over the random signs, the diagonal of
\cref{eq:ch01:spread}, though $R$ is not Haar distributed. The random signs
are needed. Without them, vectors aligned with a row of $H_d$ are concentrated
onto a single coordinate instead of being spread. Randomized Hadamard
preprocessing of this kind is the basis of the fast Johnson--Lindenstrauss
transform \cite{ailonchazelle2009}. Why a random rotation preserves distances
between many points at once, and how sharply the rotated coordinates
concentrate, is the subject of \cref{ch11}, which states the
Johnson--Lindenstrauss lemma.

\begin{example}[Spreading with a randomized Hadamard matrix]\label{ex:ch01:hadamard}
For $d=4$ Sylvester's construction gives rows $(1,1,1,1)$, $(1,-1,1,-1)$,
$(1,1,-1,-1)$ and $(1,-1,-1,1)$. Take $D_s=\diag(1,-1,1,1)$ and
$R=\tfrac12H_4D_s$. The coordinate vector $x=e_1$ maps to the first column of
$H_4$ over $2$, that is $(\tfrac12,\tfrac12,\tfrac12,\tfrac12)\T$. Every
coordinate has square $\tfrac14=\norm{x}^2/d$, and the length is still $1$.

Now take $x=\tfrac12(1,1,1,1)\T$. Without signs, $\tfrac12H_4x=(1,0,0,0)\T$. All
the energy lands on one coordinate. With the signs, $D_sx=\tfrac12(1,-1,1,1)\T$,
the dot products with the four rows of $H_4$ are $2,2,-2,2$ over $2$, and
$Rx=(\tfrac12,\tfrac12,-\tfrac12,\tfrac12)\T$, spread evenly again.

A random sign flip followed by a permutation is also orthogonal, but it maps
$e_1$ to $\pm e_j$ for some $j$. It moves energy between coordinates without
spreading it.

In dimension $8$, \cref{fig:ch01:hadamard} applies a randomized Hadamard
matrix, with signs drawn by the check script, to the spiky vector
$x=(2,0.4,0,\dots,0)\T$. The largest coordinate falls from $2$ to
$2.4/\sqrt8\approx0.849$ and every coordinate becomes $(2\pm0.4)/\sqrt8$ in
magnitude, while the length is unchanged.
\end{example}

\begin{figure}[tbp]
\centering
\begin{tikzpicture}
\begin{axis}[primer, width=0.47\linewidth, height=0.36\linewidth, ybar, bar width=6pt,
  ymin=-1.1, ymax=2.2, xtick={1,...,8}, title={\small before, $x$}]
\addplot[fill=pblue, draw=pblue] coordinates {(1,2.000) (2,0.400) (3,0.000) (4,0.000) (5,0.000) (6,0.000) (7,0.000) (8,0.000)};
\end{axis}
\end{tikzpicture}\hfill
\begin{tikzpicture}
\begin{axis}[primer, width=0.47\linewidth, height=0.36\linewidth, ybar, bar width=6pt,
  ymin=-1.1, ymax=2.2, xtick={1,...,8}, title={\small after, $Rx$}]
\addplot[fill=pgreen, draw=pgreen] coordinates {(1,0.849) (2,0.566) (3,0.849) (4,0.566) (5,0.849) (6,0.566) (7,0.849) (8,0.566)};
\end{axis}
\end{tikzpicture}
\caption{A randomized Hadamard rotation spreads a spiky vector over all eight
coordinates. The length is the same in both panels, and the largest magnitude
drops from $2$ to $0.849$, so one scalar codebook fits every coordinate.}
\label{fig:ch01:hadamard}
\end{figure}

\inproject{\nolinkurl{turboquant-pro/turboquant_pro/core.py} implements the TurboQuant
algorithm of Zandieh et al.\ \cite{zandieh2025}, which rotates and then applies a
scalar quantizer to each coordinate. The description here is of the code after
the fix recorded in \nolinkurl{turboquant-pro/CHANGELOG.md} on 2026-10-07. For
\texttt{head\_dim} up to $4096$ the rotation is \texttt{np.linalg.qr} of a
Gaussian matrix with no sign correction, so it is not exactly Haar distributed.
The changelog reports the same relative squared error with and without
Mezzadri's correction \cref{eq:ch01:haar}, $0.0444$ against $0.0453\pm0.0034$
over $20$ seeds at $d=128$, and the code comments now say so. Above $4096$ the
class \texttt{\_TwoWindowHadamard} applies a randomized Hadamard rotation of
any dimension. With $m=2^{\lfloor\log_2d\rfloor}$ it applies random signs, a
random permutation and the scaled Hadamard matrix to the first $m$
coordinates, then fresh random signs, a permutation and the Hadamard matrix
to the last $m$. Because $m>d/2$ when $d$ is not a power of two, the two
windows cover every coordinate, and when $d$ is a power of two only the first
window is used. Before the fix this branch was the sign flip plus permutation
of \cref{ex:ch01:hadamard}, which spreads no energy. On inputs with $8$ outlier
channels at $d=8192$ and $3$ bits the changelog reports relative squared error
$0.89$ for that path against $0.054$ for a spreading rotation. Gaussian inputs
did not show the difference, because they are already spread.
\nolinkurl{turboquant-pro/turboquant_pro/pgvector.py} keeps the old sign flip
plus permutation for \texttt{rotation="qr"} above $4096$ dimensions so that
stored records still decode, logs a warning, and recommends its opt-in
\texttt{rotation="hadamard"}, $R=(1/\sqrt d)H\,\diag(\text{sign})$ applied by a
fast Walsh--Hadamard transform for a power-of-two dimension. A rotation
preserves Euclidean error, which is not the acceptance criterion for a
quantizer. \Cref{ch07} explains why acceptance uses rank fidelity instead.}

\buildson{Primer~L, section~L.5. \Cref{sec:ch01:trace}.}
\usedin{\Cref{ch07}, random rotation and scalar quantization in turboquant-pro.}

% ---------------------------------------------------------------------------
\ompnexthead{The log-determinant}
\section{The log-determinant is concave, and its gradient is the inverse}\label{sec:ch01:logdet}

The determinant of a covariance is the squared volume of its unit ellipsoid,
up to a constant. Its logarithm turns products of eigenvalues into sums,
\begin{equation}\label{eq:ch01:logdet}
\log\det X=\sum_{i=1}^d\log\lambda_i(X)=\tr\log X,\qquad X\pd 0,
\end{equation}
where $\log X$ is defined by \cref{eq:ch01:funcalc}. \Cref{ch02} shows that the
mutual information between jointly Gaussian vectors is a difference of two
log-determinants, and \cref{ch03,ch04} meet a Gaussian rate--distortion problem
that maximizes one. Two calculus facts are needed for those chapters.

\begin{proposition}[Derivatives of $\log\det$ and of the inverse]\label{prop:ch01:logdetgrad}
For $X\pd0$ and a small symmetric perturbation $H$,
\begin{gather}
\log\det(X+H)=\log\det X+\tr(X^{-1}H)-\tfrac12\tr\big(X^{-1}HX^{-1}H\big)+O(\norm{H}^3), \label{eq:ch01:logdetexp}\\
\nabla_X\log\det X=X^{-1},\qquad d(X^{-1})=-X^{-1}\,(dX)\,X^{-1}. \label{eq:ch01:logdetgrad}
\end{gather}
\end{proposition}

\begin{proof}
Write $X+H=X^{1/2}(I+E)X^{1/2}$ with $E=X^{-1/2}HX^{-1/2}$. Then
$\log\det(X+H)=\log\det X+\sum_i\log(1+e_i)$, where $e_i$ are the eigenvalues
of $E$. Expanding $\log(1+e)=e-\tfrac12e^2+O(e^3)$ and using
$\sum_ie_i=\tr E=\tr(X^{-1}H)$ and $\sum_ie_i^2=\tr E^2=\tr(X^{-1}HX^{-1}H)$
gives \cref{eq:ch01:logdetexp}. The gradient is the matrix whose trace inner
product with $H$ is the first-order term. Differentiating $XX^{-1}=I$ gives
$dX\,X^{-1}+X\,d(X^{-1})=0$, which is the second identity.
\end{proof}

The second-order term $-\tfrac12\tr(E^2)$ is never positive, because $E$ is
symmetric, so $\log\det$ is concave on the positive definite matrices,
\begin{equation}\label{eq:ch01:logdetconcave}
\log\det\big(\theta X+(1-\theta)Y\big)\ \ge\ \theta\log\det X+(1-\theta)\log\det Y,\qquad 0\le\theta\le1 .
\end{equation}
It is also monotone. If $X\psd Y\pd0$, then $Y^{-1/2}XY^{-1/2}\psd I$, every
eigenvalue of it is at least $1$, and so $\log\det X\ge\log\det Y$. Concavity
alone does not give a unique maximizer, since a function that is constant on a
segment is concave. Strict concavity does, and $\log\det$ is strictly concave
because $\tr(E^2)>0$ for every symmetric $E\neq0$. So a max-det problem over a
convex feasible set has at most one maximizer, and exactly one when a
maximizer exists. Monotonicity is what makes ``smaller covariance in the
Loewner order'' mean ``no more volume''. The converse fails. $\diag(1,1)$ has
the smaller determinant of the pair $\diag(1,1)$ and $\diag(4,\tfrac12)$, yet
the two are incomparable.

\begin{example}[First and second order, and a midpoint]\label{ex:ch01:logdet}
Let $X=\begin{pmatrix}2&1\\1&1\end{pmatrix}$, with $\det X=1$ and
$X^{-1}=\begin{pmatrix}1&-1\\-1&2\end{pmatrix}$. Perturb by $H=\diag(0.1,0)$.
Exactly, $\det(X+H)=2.1-1=1.1$ and $\log1.1\approx0.09531$. The first-order term
is $\tr(X^{-1}H)=0.1$. The second-order term is
$-\tfrac12\tr(X^{-1}HX^{-1}H)=-\tfrac12(0.1)^2=-0.005$, because
$X^{-1}H$ has only one nonzero entry on its diagonal, $0.1$. So the
expansion gives $0.095$, within $0.0004$ of the exact value.

For concavity, take $X=\diag(1,4)$ and $Y=\diag(4,1)$. Each has
$\log\det=\log4\approx1.386$. Their midpoint is $\diag(2.5,2.5)$, with
$\log\det=2\log2.5\approx1.833$, larger than the average $1.386$ of the
endpoints.
\end{example}

\inproject{\nolinkurl{geometric-observation/paper/observation-theory.tex}, section
``Successive refinement across observers'', the lemma on the Gaussian optimum
in max-det form (label \texttt{lem:maxdet}), writes the consumer-relative rate as
$\tfrac12\log\det\Sigma_x-\tfrac12\log\det\Sigma^\star$ with $\Sigma^\star$
maximizing $\log\det\Sigma$ subject to $0\preceq\Sigma\preceq\Sigma_x$ and
$\tr(P\Sigma)\le D$. Its uniqueness argument is the strict concavity in
\cref{eq:ch01:logdetconcave}.}

\buildson{Primer~L, section~L.7 (determinant as product of eigenvalues) and
section~L.9 (gradients). \Cref{sec:ch01:spectral}.}
\usedin{\Cref{ch02}, Gaussian mutual information. \Cref{ch03}, Gaussian
rate--distortion. \Cref{ch04}, the max-det program and its KKT conditions.}

% ---------------------------------------------------------------------------
% ---------------------------------------------------------------------------
\ompnexthead{The tools together}
\section{The tools combine to price a direction the coder did not read}\label{sec:ch01:together}

This last section puts five tools of the chapter into one calculation from
the programme. A coder designs its error for an \emph{estimated} read
operator $\hat P$ while the consumer's true operator is $P$. Suppose the coder
codes perfectly every direction that $\hat P$ reads and leaves the rest alone.
How much distortion does the true consumer still see? The answer needs a
whitened operator (\cref{sec:ch01:whiten}), a kernel and its orthogonal
projector (\cref{sec:ch01:proj}), a trace (\cref{sec:ch01:trace}) and, for
interpretation, a Schur complement (\cref{sec:ch01:schur}).

Work in whitened coordinates $z=\Sigma_x^{-1/2}x$, where the source has
covariance $I$. The coder reads the directions outside
$\ker(\Sigma_x^{1/2}\hat P\Sigma_x^{1/2})$ and leaves the whitened kernel
uncoded. Let $\Pi$ be the orthogonal projector onto that kernel. The whitened
error is then $\Pi z$, with covariance $\Pi I\Pi=\Pi$, and by
\cref{eq:ch01:whitenedop} the true consumer sees
\begin{equation}\label{eq:ch01:floor}
D_{\mathrm{floor}}=\tr\big(\tilde P\,\Pi\big),\qquad \tilde P=\Sigma_x^{1/2}P\,\Sigma_x^{1/2},\qquad
\Pi=\text{projector onto }\ker\big(\Sigma_x^{1/2}\hat P\,\Sigma_x^{1/2}\big).
\end{equation}
By \cref{eq:ch01:tracepos} the floor is positive exactly when $\tilde P\Pi\neq0$,
that is, when the true consumer reads some whitened direction the coder left
out. The whitening matters. Using the projector onto $\ker\hat P$ in source
coordinates ignores that coding the read directions also reduces the error in
whatever is correlated with them.

\begin{example}[An omitted direction that is partly recovered]\label{ex:ch01:floor}
Let $\Sigma_x=\begin{pmatrix}2&1\\1&1\end{pmatrix}$. The coder believes the
consumer reads $x_1$, so $\hat P=e_1e_1\T$. The true consumer reads $x_2$, so
$P=e_2e_2\T$. Then $\Sigma_x^{1/2}\hat P\Sigma_x^{1/2}=s_1s_1\T$ with
$s_1=\Sigma_x^{1/2}e_1$, so its kernel is the line orthogonal to $s_1$ and
$\Pi=I-s_1s_1\T/\norm{s_1}^2$, with $\norm{s_1}^2=e_1\T\Sigma_xe_1=2$. Using
the cyclic rule,
\[
D_{\mathrm{floor}}=\tr\tilde P-\frac{s_1\T\tilde Ps_1}{2}
=e_2\T\Sigma_xe_2-\frac{(e_1\T\Sigma_xe_2)^2}{2}=1-\tfrac12=\tfrac12 .
\]
This is the Schur complement $\Sigma_{22}-\Sigma_{21}\Sigma_{11}^{-1}\Sigma_{12}$
of \cref{eq:ch01:schur}. For a Gaussian source \cref{ch02} identifies it as the
variance of $x_2$ left after $x_1$ is known exactly. For any other source it is
the error variance of the best \emph{linear} predictor of $x_2$ from $x_1$, and
a nonlinear decoder can do better. Coding $x_1$ perfectly has removed half of
the variance of $x_2$. The naive projector onto $\ker\hat P=\operatorname{span}(e_2)$,
applied to $\tilde P$, gives $\big((\Sigma_x^{1/2})_{22}\big)^2$. For a
$2\times2$ positive definite matrix the square root is
$(\Sigma_x+\sqrt{\det\Sigma_x}\,I)/\sqrt{\tr\Sigma_x+2\sqrt{\det\Sigma_x}}$,
here $\tfrac1{\sqrt5}\begin{pmatrix}3&1\\1&2\end{pmatrix}$, so the naive number
is $\tfrac45$. It overstates the floor of $\tfrac12$.
\end{example}

\inproject{\nolinkurl{geometric-observation/paper/observation-theory.tex}, appendix ``The
price of a misidentified observer'', theorem ``Omission floor'', defines
$D_{\mathrm{floor}}:=\tr(\tilde P\Pi)$ with $\Pi$ projecting onto the whitened
kernel, proves that every allocation in the $\hat P$-water-filling family has
true distortion at least $D_{\mathrm{floor}}$, and remarks that ``the naive
$\ker\hat P$ overstates the floor when $\Sigma_x\ne I$''.
Example~\ref{ex:ch01:floor} is a two-dimensional instance of that remark. The
overstatement is not automatic. With $\Sigma_x=\diag(2,1)$ and $\hat P=e_1e_1\T$
the whitened kernel $\Sigma_x^{-1/2}\operatorname{span}(e_2)$ is
$\operatorname{span}(e_2)$ itself, so the naive projector gives the floor exactly.}

\buildson{\Cref{sec:ch01:proj,sec:ch01:trace,sec:ch01:whiten,sec:ch01:schur}.}
\usedin{\Cref{ch02}, where the same Schur complement is a conditional variance.
\Cref{ch06}, the commission tax and the omission floor.}

\let\sectionmark\ompSavedSectionmark
\section*{Exercises}

\subsection*{Warm-up}

\begin{exercise}\label{exr:ch01:rank}
Show that the trace of an orthogonal projector equals its rank. Check it on the
projector of Example~\ref{ex:ch01:gproj} in $\R^3$.
\end{exercise}

\begin{exercise}\label{exr:ch01:oblique}
Is $\Pi=\begin{pmatrix}1&1\\0&0\end{pmatrix}$ a projector? Is it an orthogonal
projector? Find its range and kernel and the angle between them.
\end{exercise}

\begin{exercise}\label{exr:ch01:incomparable}
For $A=\diag(3,1)$ and $B=\begin{pmatrix}2&1\\1&2\end{pmatrix}$, decide whether
$A\psd B$, $B\psd A$ or neither.
\end{exercise}

\begin{exercise}\label{exr:ch01:cyclic}
Find $2\times2$ matrices with $\tr(ABC)\neq\tr(ACB)$.
\end{exercise}

\subsection*{By hand}

\begin{exercise}\label{exr:ch01:whiten}
Let $\Sigma_x=\diag(4,1)$ and $P=gg\T$ with $g=(1,-1)\T$. Compute the whitened
operator $\tilde P$, its nonzero eigenvalue and eigenvector, and $\tr(P\Sigma_x)$.
\end{exercise}

\begin{exercise}\label{exr:ch01:pencil}
Find the generalized eigenvalues and eigenvectors of the pencil $(A,B)$ with
$A=\diag(1,0)$ and $B=\begin{pmatrix}2&1\\1&1\end{pmatrix}$. Which direction
maximizes $x\T Ax/x\T Bx$, and what is the maximum?
\end{exercise}

\begin{exercise}\label{exr:ch01:schur}
For $M=\begin{pmatrix}2&1&1\\1&2&1\\1&1&2\end{pmatrix}$, compute the Schur
complement of the top-left entry, decide whether $M\pd0$, and compute $\det M$
from \cref{eq:ch01:schurdet}.
\end{exercise}

\begin{exercise}\label{exr:ch01:signperm}
Before its 2026-10-07 fix, turboquant-pro rotated vectors of more than $4096$
dimensions by a sign flip followed by a permutation. Show that such a map sends
every coordinate vector to a signed coordinate vector, and that
$\tfrac1{\sqrt d}H_dD_s$ sends every coordinate vector to a vector whose
coordinates all have magnitude $1/\sqrt d$. Then show that the two Hadamard
windows $[0,m)$ and $[d-m,d)$ of the current rotation, with
$m=2^{\lfloor\log_2d\rfloor}$, together cover every coordinate.
\end{exercise}

\begin{exercise}\label{exr:ch01:logdet}
Let $X=\diag(2,4)$ and $H=\begin{pmatrix}0&0.1\\0.1&0\end{pmatrix}$. Compute
$\log\det(X+H)-\log\det X$ exactly and compare with the first- and second-order
terms of \cref{eq:ch01:logdetexp}.
\end{exercise}

\subsection*{Programming}

\begin{exercise}\label{exr:ch01:mc}
Draw $10^5$ errors $\delta=\mu+\Sigma^{1/2}z$ with $z\sim\Normal(0,I)$, using
$\Sigma=\Sigma_a$, $\mu=(1,0)\T$ and $P=\PC$ from Example~\ref{ex:ch01:trace}. Compare the
sample mean of $\delta\T P\delta$ with $\tr(P\Sigma)+\mu\T P\mu=10$.
\end{exercise}

\begin{exercise}\label{exr:ch01:haar}
Generate $10^4$ random $2\times2$ orthogonal matrices with \texttt{np.linalg.qr}
of a Gaussian matrix, with and without the sign correction of \cref{eq:ch01:haar}.
Record the determinant and the rotation angle of each. Compare the two
histograms with the uniform distribution that Haar measure gives, and check
\cref{eq:ch01:spread} for $x=(3,4)\T$.
\end{exercise}

\begin{exercise}\label{exr:ch01:kyfan}
For a random symmetric $6\times6$ matrix and $k=2$, evaluate $\tr(U\T AU)$ for
$10^4$ random $U$ with orthonormal columns and confirm that none exceeds
$\lambda_1+\lambda_2$, which the top two eigenvectors attain.
\end{exercise}

% checks/ch01_examples.py on Atlas: 179/179 PASS
